\chapter{Conclusions}
\label{cap:conclusions}

This thesis aimed to address the limitations of existing high-resolution snow products available for Northern Italy by developing a modelling 
framework based solely on historical temperature and precipitation measurements. This model was calibrated against in situ snow depth records by 
maximizing the agreement between simulated and observed snow covered periods. The calibration focused on the temperature threshold for phase 
partitioning, which converged to $T_\tau \approx \SI{1.1}{\degreeCelsius}$, while the melt related parameters were fixed to literature values 
(see Chapter \ref{cap:swe_calibration}). Additionally, a simulated annealing scheme was employed to investigate whether jointly varying $T_\tau$, 
$\text{DDF}_{ave}$ and $\text{DDF}_{amp}$ could further enhance model performance. However, no alternative parameter configuration capable of 
yielding substantial improvements was identified. Overall, the proposed approach enabled the reconstruction of snow water equivalent dynamics over 
the 1951--2023 period, providing a robust basis for evaluating the impacts of climate change on seasonal snowpack across the IH-GAR region. 

The results reveal a widespread and significant decline in snow cover across most of the study area, strongly linked to ongoing global warming, as 
the mean temperature over the IH-GAR region has increased by nearly \SI{2}{\degreeCelsius} since 1985 (see Chapter \ref{cap:swe_model}). This 
temperature rise has a twofold negative impact on snow resources: it reduces the amount of precipitation occurring as snowfall and simultaneously
enhances snowmelt processes. The magnitude of these changes becomes evident when examining the annual peak SWE time series, which is showing 
a statistically significant negative trend of -\SI{0.567}{\giga\meter\cubed} per decade over the 1951--2023 period. Such a decline corresponds, 
on average, to a reduction of approximately 35-40\% in peak SWE values over a time span of just 70 years. Furthermore, elevation bands below 
\SI{1500}{\meter} a.s.l. emerge as the most vulnerable to climate change, as winter temperatures hover around \SI{0}{\degreeCelsius}, so that even 
a slight warming leads to marked shifts in snow accumulation and melt patterns.




\section{Future developments}

Through the analyses presented in this thesis, it was possible to evaluate the influence of climate change on the seasonal snow cover across IH-GAR region 
over the 1951--2023 period. However, some limitations highlighted by the calibration procedure suggest that the modelling approach could be further 
refined to ensure a more accurate representation of snowpack physical processes. In particular, three potential improvements to the snow model can be 
identified:
\begin{itemize}[label=$\blacktriangleright$]
    \item the cold content of snow cover may be monitored throughout the snow cycle to account for the time lag between the onset of favourable 
    conditions for snowmelt and the beginning of actual runoff. As a first parametrization attempt, an approach similar to that of IT-SNOW could be 
    adopted, based on the mean air temperature computed over a time window of several days preceding the potential melt event. Additionally, the width 
    of the time window could be adjusted as a function of snow water equivalent in an attempt to consider into the modelling chain the longer time 
    required for snowpacks to ripen. 
    \item refreezing processes could also be taken into account, particularly in deeper snowpacks, where subfreezing layers may still persist even when 
    the upper layers are undergoing wet metamorphism and water percolation. 
    \item an alternative parametrization of precipitation phase partitioning may be considered, with $T_\tau$ as a time dependent parameter defined 
    either as a periodically varying function over the entire hydrological year or as a step function representing different stages of the snow season. 
    Such an approach could account for seasonal variations in atmospheric conditions influencing precipitation phase. In particular, in Alpine regions, 
    the vertical atmospheric temperature gradient is typically less pronounced during winter than in spring. This implies that, to reach a location 
    where the near-surface air temperature exceeds \SI{0}{\degreeCelsius}, snowflakes in winter must fall through a considerably deeper above-freezing 
    atmospheric layer rather than in spring, thereby increasing the likelihood of rainfall to occur. This phenomenon is responsible for the increased 
    frequency of snowfall events at 2/3 \SI{}{\degreeCelsius} during spring, while almost none are recorded during winter.
\end{itemize}
Moreover, the modelling framework developed in this thesis could be employed to investigate future changes in snow cover under different emission 
scenarios. The required meteorological forcing data are currently being downscaled thanks to a collaboration between the Department of Environmental 
Science and Policy of the University of Milan and ISAC/CNR.
